\section{Discussion}
\label{sec:discussion}

This study asked whether a geometry-based method could generate synthetic profiles from a
very small longitudinal cohort while preserving statistical and modeled biological
relationships. Median feature MRE was \rone{1.2\%}, which the MVN baseline also achieved.
Stochastic attention also retained cross-visit covariance, whereas MVN introduced variance in
194 null dimensions. Conditioned generation preserved the subgroup patterns represented by
three PCOS and five Developed~PE patients. Finally, \rboth{a separately specified,
generator-blind coagulation model calibrated on the same cohort placed real and synthetic
profiles in the same factor-to-thrombin-generation envelope}. Cloud overlap was 83--92\%, and
the Kolmogorov--Smirnov tests did not detect differences under these conditions ($p>0.35$ for
all five TGA features). A model calibrated on synthetic profiles predicted held-out real TGA
measurements as well as one calibrated on real profiles. \rone{Direct sampling from the same
target distribution reproduced the Langevin sampler's fidelity and membership-inference
results.} Together, these findings provide \rone{proof of concept for using synthetic profiles
in the modeling tasks tested here} with very small longitudinal cohorts.

Dense simulations from very small longitudinal datasets could support maternal health
research. Rare pregnancy complications such as PCOS and preeclampsia are difficult to study
because assembling large, well-characterized longitudinal cohorts requires years of
recruitment across multiple clinical sites. Our results support hypothesis generation,
computational modeling, and workflow testing with dense simulations of these small cohorts.
\rtwo{Related work has predicted TGA responses from coagulation factors and searched for
factor changes that produce a target response~\citep{Ghetmiri2021}. Although we did not
evaluate either task with profiles generated by stochastic attention, these applications
offer a clear direction for extending the present work.}

\rboth{Understanding why stochastic attention worked in this setting required examining the role of the
PCA representation as well as the generator. When all methods used
the same $18$-dimensional subspace, the fitted two-component mixture and copula achieved
marginal errors and mechanistic overlaps similar to those of stochastic attention. These results were therefore
not unique to stochastic attention, although the comparison did not quantify the contribution
of PCA itself. Empirical correlation error and novelty differed across generators. Stochastic
attention used the patient profiles in the PCA subspace as its memory matrix and avoided
estimating a full $216$-dimensional covariance from $23$ patients. Its sampler still required
a choice of inverse temperature.} A natural concern is whether the $\beta^*$
selection procedure, originally developed for protein sequence
generation~\citep{Varner2026SAProtein}, transfers to clinical coagulation data. The entropy
transition rule selected the point of greatest downward curvature in attention entropy, a
geometric property of $K$ patterns in $d$ dimensions rather than a property of what the
features represented. For this dataset it gave $\beta^* \approx 2.94$, with the random-pattern
scale $\sqrt{d} \approx 4.24$ serving as a theoretical reference. A sensitivity analysis
showed gradual changes in generation quality across $\beta/\beta^* \in [0.1, 3.0]$
(\supptable{stab:beta-sweep}).
Similarly, varying the PCA variance retention threshold from 85\% to 99\%
($d_{\text{PCA}} = 13$ to 21) had minimal impact on generation quality. Median MRE
remained between 1.2\% and 1.8\% across all thresholds (\supptable{stab:pca-sensitivity}),
so the 95\% threshold was not a critical design choice. \rboth{Together, these sensitivity
analyses showed that performance was not tied to a narrow choice of inverse temperature or
PCA threshold within the ranges tested.}

Conditional generation raised a different question: how much diversity remained after the
selected subgroup was upweighted? Achieving
80\% of finite-mixture component probability on the PCOS subgroup required
$\rho \approx 26.7$, reducing the participation ratio to
$K_{\text{eff}} \approx 4.6$ (multiplicity parameters
for all three subgroups in \supptable{stab:multiplicity}),
which indicated that mixture-component probability was concentrated around a small number
of stored profiles across repeated draws. The subgroup-specific signal still rested on only
three real PCOS profiles, and the lower-weight background profiles pulled generated means
toward the population average; $K_{\text{eff}}$ did not represent an independent patient
count. Magnitudes were drawn from all $23$ empirical PCA-space norms, so conditioning acted
through the weighted directions rather than a subgroup-specific magnitude distribution.
Thus, weighting concentrated sampling toward the subgroup directions while retaining some
influence from the full cohort.
\rboth{Related studies have applied the same geometry-based operating-point rule
and multiplicity weighting to protein sequence generation and conditional
steering~\citep{Varner2026SAProtein,Varner2026Multiplicity,Alswaidan2026SA}.}

Among longitudinal clinical generators, VAMBN-MT provides a useful comparison in a
larger-data setting~\citep{Kuhnel2024}. It combines
expert-defined modules, Bayesian-network constraints, and LSTM-augmented variational
autoencoders. The authors evaluated it on the DONALD nutritional cohort and the Alzheimer's
Disease Neuroimaging Initiative. DONALD included $1{,}312$ participants and 33 longitudinal
variables, compared with 23 patients and 216 features here. Its best variant had cross-visit
correlation error near 0.70, and weaker time trends were not reproduced at the tested sample
sizes. The methods address different regimes. The VAMBN-MT model uses a trained modular architecture
for moderate cohorts. Stochastic attention uses PCA geometry when $n<p$, applies a generic
generator-blind mechanistic check instead of domain-specific rules, and uses multiplicity
weights for conditional generation. This was a regime comparison, not a head-to-head test.

The mechanistic analysis addressed a different validation question: whether a model of the
underlying biology processed real and synthetic profiles similarly. We processed both groups
with \rboth{the same separately specified model, which was blind to the generator but
calibrated on the real cohort}. A model calibrated on synthetic profiles also
predicted held-out real TGA measurements as well as one calibrated on real profiles (error
ratio $0.94$, with 11 random restarts). The same comparison can be used wherever a validated
mechanistic model exists, including pharmacokinetic~\citep{MouldUpton2013},
physiological~\citep{Chase2011}, tumor-growth~\citep{Ribba2012}, and
metabolic~\citep{Yizhak2015,Shan2018} models. The relevant test is whether the model places
real and synthetic profiles in the same predicted envelope.

\rone{The subgroup analysis focused on average differences between groups, but it did not
show whether the conditioned cohorts preserved clinically important extremes. In the
descriptive bootstrap Mann--Whitney diagnostic, the test detected no difference in fewer
than 90\% of replicates for four of the 24 feature--condition pairs: FIX and plasminogen in
the PCOS and Developed~PE groups. For these pairs, the relative differences between the real
and synthetic subgroup means were 9--15\%, and values varied widely among the real patients.
This diagnostic had limited power and was not an equivalence test. The lower fractions could
therefore reflect shifts in the means, a loss of extreme values, or simply the small samples.
Truncating the PCA representation and drawing profile magnitudes from only $23$ observed
values could also make the generated extremes too narrow. If so, a synthetic cohort could
not show how often severe states occur or how severe they become. Profiles outside the real cohort's
convex hull did not show that the true extremes were reproduced, and the primary mechanistic
cloud-overlap measure considered only the real 5th--95th percentile range. We therefore do
not recommend using these cohorts to estimate extreme risk without additional real data.}
\rboth{The small source cohort caused other limits. Hull membership had little resolution,
and the conditioned cohorts still relied on only three PCOS or five Developed~PE patients.
The modest calibration improvement was consistent with drawing more profiles from the same
patient-derived distribution, not with adding biological information. The generated profiles
supported simulation and model testing, but they did not add clinical evidence beyond the
source cohort.}

\rone{The pipeline required every patient to have the same assays at the same aligned visits.
Concatenating the visits preserved relationships across visits, but the pipeline could not
directly use irregular visit times, missing assays or visits, or different numbers of visits.
Imputation, padding, or a sequence-aware method could support these records, but each would
add assumptions that require validation in a small cohort.}
\rtwo{PCA added another limit because it captured only linear patterns. In the S-curve
benchmark, generated profiles moved farther from the known curve as curvature increased.
At every nonzero curvature, stochastic attention remained closer to the curve on average
than the Gaussian generator, whereas the Gaussian generator recovered covariance more
accurately. Neither generator was best on both measures. A nonlinear representation could
better follow curved or branching patterns, but it would add choices about fitting the model
and converting the generated values back to patient measurements. Those choices are
difficult to make from $23$ patients. The linear representation was practical for this
cohort, but the results did not show that clinical trajectories are generally linear.}
\rone{Cohort size also affects computational cost. The reported Langevin sampler required
work proportional to $TKd$ for each profile, where $T$ was the number of Langevin steps,
$K$ was the number of stored profiles, and $d$ was the dimension of each profile in the
sampling space. The exact mixture sampler removed the $T$ sequential updates: its component
weights could be prepared once, after which each profile
required a component draw, a $d$-dimensional Gaussian draw, and conversion back to the
original variables. The simulation advantage over the sample-covariance Gaussian occurred
mainly in the $n<p$ settings, and we did not test cohorts larger than $23$ patients.}
\rboth{The results therefore did not establish performance on irregular, incomplete, or
strongly nonlinear records or in larger clinical cohorts.}

This study had two broad limitations: the absence of independent clinical data and the
restricted scope of the validation analyses. First, the clinical analysis used one
longitudinal cohort of $23$ patients. \rboth{The known-covariance simulations tested selected
properties under controlled conditions, and related protein studies applied the sampling
framework in another domain, but neither established that the clinical findings would
generalize to a different patient population. We did not test the method in a second clinical
cohort, which is needed to assess that generalizability.} Second, the validation analyses
answered narrower questions than full biological or clinical validity. \rboth{The BZ2012
model was calibrated on the same cohort used for generation. Applying one fixed,
generator-blind model to real and synthetic profiles therefore tested consistency with that
model, not independent biological validity.}
The fitted rate constants also required a 50-fold reduction in intrinsic tenase
$k_{\text{cat}}$ and a 15-fold
increase in extrinsic tenase $k_{\text{cat}}$ relative to published values
(\supptable{stab:ode-calibration}), likely because the published measurements and the
plasma-based assay used different experimental conditions. These fitted values should not be
read as new biochemical estimates. \rone{The biological scope was also incomplete. The
mechanistic test covered thrombin generation but not the fibrinolytic or viscoelastic
features. The membership-inference analysis answered a separate, narrow question: whether an
already-known de-identified profile had been stored in the memory matrix. It distinguished
stored from held-out profiles, but it did not identify a participant or reconstruct an
unknown record.} Finally, we did not test clinical outcomes. Thus, the results
did not establish biological fidelity across the full profile or clinical utility. An
independently calibrated model, models covering the other feature classes, and prospective
outcome validation could address these gaps. The results establish proof of concept, not
generalizability.
